\section{Разработка низкоуровневого ПО для МК}
\label{sec:firmware}

Разработка низкоуровневого программного обеспечения (прошивки) была направлена на создание базового функционала системы: сбор данных с датчиков, их первичную обработку и организацию надёжного двустороннего обмена данными с персональным компьютером через интерфейс USB. Конкретными задачами являлись инициализация аппаратных модулей МК, реализация драйверов для гироскопа L3GD20 и комбинированного датчика LSM303DLHC, организация циклического опроса сенсоров, а также создание протокола для передачи телеметрии и приёма управляющих команд.

Разработка велась на языках С и С++, что соответствует отраслевым стандартам для встраиваемых систем благодаря высокой эффективности и низкоуровневому контролю. В качестве основы для работы с периферией микроконтроллера использовалась библиотека стандартной периферии (Standard Periphery Library, SPL) от STMicroelectronics. Данная библиотека предоставляет структурированный набор функций для конфигурации всех встроенных модулей МК, что ускоряет разработку, сохраняя прозрачность управления аппаратными ресурсами.